% !TEX root = ../main.tex
\section{Proofs of the Upper Bounds}
\label{app:upper-bounds}

\subsection{Game stationarity and KKT points}
\label{app:kkt-gs}

We next relate Definition~\ref{def:eps-kkt} to game stationarity (GS) for the
minimax reformulation~Problem~\eqref{eq:ncl} of Problem~\eqref{prob:main}. We use the
following GS definition from \cite[Definition~3.1(b)]{li2025nonsmooth} and
\cite[(2.10)--(2.11)]{zhang2026matching}.
\begin{definition}[\(\epsilon\)-GS point]
\label{def:eps-gs}
Let \(\epsilon>0\). A pair \((\bx,\bq)\in\RR^n\times\dom\partial h^*\) is an
\(\epsilon\)-game stationary point (\(\epsilon\)-GS) of~Problem~\eqref{eq:ncl} if
\[
    \calG(\bx,\bq)
    :=
    \max\Bigl\{
        \bigl\|\nabla f(\bx)+J_{\calA}(\bx)^\top\bq\bigr\|,\;
        \dist\bigl(\calA(\bx),\partial h^*(\bq)\bigr)
    \Bigr\}
    \le\epsilon .
\]
\end{definition}
We now relate $\epsilon$-KKT points of Problem~\eqref{prob:main}
to $\epsilon$-GS points of its Fenchel saddle reformulation~Problem~\eqref{eq:ncl}. 
\begin{proposition}[KKT points of Problem~\eqref{prob:main} and GS points of~Problem~\eqref{eq:ncl}]
\label{prop:kkt-gs}
Suppose that Assumption~\ref{ass:h} holds, and let
\((\bx,\bq)\in\RR^n\times\dom\partial h^*\). Then the following measurement relationship holds: 
\[
    \calG(\bx,\bq)
    =
    \min_{\by\in\partial h^*(\bq)}
    \mathcal{R}_{\mathrm{KKT}}(\bx,\by,\bq),
\]
with minimum attained at
\(\hat\by=\Pi_{\partial h^*(\bq)}\bigl(\calA(\bx)\bigr)\). Consequently, for
every \(\epsilon>0\):
\begin{enumerate}[label=(\alph*)]
\item if \((\bx,\by,\bq)\) is an \(\epsilon\)-KKT point
  of~Problem~\eqref{prob:main}, then \((\bx,\bq)\) is an \(\epsilon\)-GS point
  of~Problem~\eqref{eq:ncl};
\item if \((\bx,\bq)\) is an \(\epsilon\)-GS point of~Problem~\eqref{eq:ncl}, then
  \((\bx,\hat\by,\bq)\) is an \(\epsilon\)-KKT point of~Problem~\eqref{prob:main}.
\end{enumerate}
\end{proposition}

\begin{proof}
Since \(h\) is proper, closed, and convex,
\(\bq\in\partial h(\by)\) if and only if \(\by\in\partial h^*(\bq)\)
\cite[Theorem~23.5]{rockafellar1970convex}. The first term of
\(\mathcal{R}_{\mathrm{KKT}}\) does not depend on \(\by\). The set
\(\partial h^*(\bq)\) is nonempty, closed, and convex, so
\(\min_{\by\in\partial h^*(\bq)}\|\calA(\bx)-\by\|
=\dist(\calA(\bx),\partial h^*(\bq))\), attained at \(\hat\by\).
Statements~(a) and~(b) follow.
\end{proof}

\subsection{Proof of Theorem~\ref{thm:upper-beta1}}
\label{app:upper-beta1}

\subsubsection{Proof of the iteration complexity bound}

Throughout this appendix, let
\[
L:=\sqrt2\,\max\Bigl\{L_{\calD,\bx},\ G_{\calA},\ \frac1\rho\Bigr\},
\qquad
\bar\gamma:=\frac1{4L},
\qquad
\bar\tau:=\frac1{36L}\Bigl(\frac\eta\gamma\Bigr)^2,
\]
and let \(\mathcal Y:=\{\blmbd\in\RR^m:\|\blmbd\|\le L_h+1\}\). We write
\([\mathrm{ZXSL}]\) for~\cite{zhang2020smoothed}; equation and lemma numbers
refer to its arXiv version~2010.15768v3.

\paragraph{Step 1: reduction to Smoothed GDA}
Since \(\beta=1\) and \(\bmu^0=\blmbd^0=\zero\), induction gives
\(\bmu^k=\blmbd^k\) for all \(k\). Hence the gradient of the dual regularizer
vanishes at every dual step, and the updates~\eqref{eq:splalm-updates} read
\[
\bx^{k+1}=\bx^k-\eta\Bigl[\nabla_{\bx}\calD_\rho(\bx^k,\blmbd^k)
  +\frac1\gamma(\bx^k-\bz^k)\Bigr],
\qquad
\blmbd^{k+1}=\blmbd^k+\tau\nabla_{\blmbd}\calD_\rho(\bx^{k+1},\blmbd^k),
\]
together with \(\bz^{k+1}=\bz^k+\alpha(\bx^{k+1}-\bz^k)\). None of these
updates involves \(s\). By~\eqref{eq:prox-multiplier},
\(\nabla_{\blmbd}\calD_\rho(\bx,\blmbd)=\calA(\bx)-\by(\bx,\blmbd)
=(\bq(\bx,\blmbd)-\blmbd)/\rho\), so
\[
\blmbd^{k+1}=\Bigl(1-\frac\tau\rho\Bigr)\blmbd^k+\frac\tau\rho\,\bq^k,
\qquad \|\bq^k\|\le L_h,
\]
where the bound holds because \(\bq^k\in\partial h(\by^{k+1})\) and \(h\) is
\(L_h\)-Lipschitz. Since \(L\ge1/\rho\), we have
\(\tau\le\bar\tau\le1/(144L)\le\rho\), and induction gives
\(\|\blmbd^k\|\le L_h\) for all \(k\). Hence every \(\blmbd^k\) lies in the
interior of \(\mathcal Y\), and the projection onto \(\mathcal Y\) does not
change the iterates. Consequently, \(\{(\bx^k,\blmbd^k,\bz^k)\}\) is the
sequence generated by Smoothed GDA~\cite[Algorithm~2]{zhang2020smoothed}
applied to
\(\min_{\bx\in\RR^n}\max_{\blmbd\in\mathcal Y}\calD_\rho(\bx,\blmbd)\) from
zero, with smoothing parameter \(p=1/\gamma\), primal stepsize \(\eta\), dual
stepsize \(\tau\), and averaging parameter \(\alpha\).

\paragraph{Step 2: assumptions of~\cite[Theorem~3.4]{zhang2020smoothed}}
The proof of Proposition~\ref{prop:curvature-smoothness} gives
\[
\begin{aligned}
\|\nabla_{\bx}\calD_\rho(\bx,\blmbd)-\nabla_{\bx}\calD_\rho(\bx',\blmbd')\|
&\le L_{\calD,\bx}\|\bx-\bx'\|+G_{\calA}\|\blmbd-\blmbd'\|,\\
\|\nabla_{\blmbd}\calD_\rho(\bx,\blmbd)-\nabla_{\blmbd}\calD_\rho(\bx',\blmbd')\|
&\le G_{\calA}\|\bx-\bx'\|+\frac1\rho\|\blmbd-\blmbd'\|.
\end{aligned}
\]
By the Cauchy--Schwarz inequality, both gradients are \(L\)-Lipschitz on
\(\RR^n\times\RR^m\). The set \(\mathcal Y\) is convex and compact,
\(\RR^n\) is closed and convex, and \(\calD_\rho(\bx,\cdot)\) is concave. By
Lemma~\ref{lem:effective-penalty},
\[
\calD_\rho(\bx,\blmbd)
=\max_{\bq}\Bigl\{\calL(\bx,\bq)-\frac1{2\rho}\|\bq-\blmbd\|^2\Bigr\}
\le\max_{\bq}\calL(\bx,\bq)=F(\bx),
\]
with equality for \(\blmbd\in\partial h(\calA(\bx))\). Every such
\(\blmbd\) satisfies \(\|\blmbd\|\le L_h\) and hence lies in \(\mathcal Y\).
Therefore \(\max_{\blmbd\in\mathcal Y}\calD_\rho(\bx,\blmbd)=F(\bx)\), which
is bounded below by \(\inf F>-\infty\). Thus Assumption~3.3 of
\([\mathrm{ZXSL}]\) holds with gradient Lipschitz constant \(L\).

\paragraph{Step 3: parameter conditions}
We verify condition~(3.3) of \([\mathrm{ZXSL}]\). First,
\(\gamma\le\bar\gamma\) gives \(p=1/\gamma\ge4L>3L\). Second, since
\(\gamma L\le1/4\),
\(\gamma/(1+\gamma L)\ge4\gamma/5>\gamma/2\ge\eta\), hence
\(\eta<1/(p+L)\). Third, put \(a:=\eta(p-L)\). Then
\(3\eta/(4\gamma)\le a<\eta p\le1/2\), so
\[
\frac{\eta^2(p-L)^2}{4L(1+\eta(p-L))^2}
>\frac{a^2}{16L}
\ge\frac9{256L}\Bigl(\frac\eta\gamma\Bigr)^2
>\bar\tau\ge\tau,
\]
and \(\bar\tau\le1/(144L)<1/(11L)\). Finally, the averaging parameter
\(\alpha=c\epsilon^2\le c\) satisfies
\(\alpha\le\min\{1/36,(p-L)^2/(384p(p+L)^2)\}\) by the choice of \(c\) in
Step~5.

\paragraph{Step 4: initial gap}
Let \(K(\bx,\bz;\blmbd):=\calD_\rho(\bx,\blmbd)+\frac1{2\gamma}\|\bx-\bz\|^2\),
\(d(\blmbd,\bz):=\min_{\bx}K(\bx,\bz;\blmbd)\), and
\(P(\bz):=\min_{\bx}\max_{\blmbd\in\mathcal Y}K(\bx,\bz;\blmbd)\). The
potential function of~\cite[(4.1)]{zhang2020smoothed} at the zero
initialization is \(\varphi^0=K(\zero,\zero;\zero)-2d(\zero,\zero)+2P(\zero)\).
By Proposition~\ref{prop:curvature-smoothness}, \(K(\cdot,\zero;\zero)\) is
\(\mu_\gamma\)-strongly convex, and
\(\|\nabla_{\bx}\calD_\rho(\zero,\zero)\|
=\|\nabla f(\zero)+J_{\calA}(\zero)^\top\bq(\zero,\zero)\|\le G_f+G_{\calA}L_h\).
Hence
\[
d(\zero,\zero)\ge\calD_\rho(\zero,\zero)-\frac{(G_f+G_{\calA}L_h)^2}{2\mu_\gamma}.
\]
By Step~2, \(P(\zero)=\min_{\bx}\{F(\bx)+\frac1{2\gamma}\|\bx\|^2\}\le F(\zero)\).
Moreover,
\(\calD_\rho(\zero,\zero)=f(\zero)+\min_{\by}\{h(\by)+\frac\rho2\|\calA(\zero)-\by\|^2\}\),
and the Lipschitz continuity of \(h\) gives
\(0\le F(\zero)-\calD_\rho(\zero,\zero)\le L_h^2/(2\rho)\). Combining these
bounds,
\[
\varphi^0-\inf F\le\hat\Delta
:=\Delta+\frac{L_h^2}{2\rho}+\frac{(G_f+G_{\calA}L_h)^2}{\mu_\gamma},
\]
where \(\mu_\gamma\ge1/\gamma-L\ge3/(4\gamma)>0\).

\paragraph{Step 5: rate}
For \(t\ge0\), set
\[
S_{t+1}:=\max\bigl\{\|\nabla_{\bx}\calD_\rho(\bx^{t+1},\blmbd^{t+1})\|,\
\|\nabla_{\blmbd}\calD_\rho(\bx^{t+1},\blmbd^{t+1})\|\bigr\}.
\]
Since \(\blmbd^{t+1}\) lies in the interior of \(\mathcal Y\), the normal
cone of \(\mathcal Y\) at \(\blmbd^{t+1}\) is \(\{\zero\}\), and
\((\bx^{t+1},\blmbd^{t+1})\) is an \(S_{t+1}\)-solution in the sense
of~\cite[Definition~3.1]{zhang2020smoothed}. Fix \(T\ge1\). The proof
of~\cite[Theorem~3.4]{zhang2020smoothed} distinguishes the cases (B.52) and
(B.53) and yields some \(t\in\{0,\dots,T-1\}\) with either
\[
\text{(i)}\quad S_{t+1}\le\lambda_1\max\{\alpha,\sqrt\alpha\}
\qquad\text{or}\qquad
\text{(ii)}\quad S_{t+1}\le\lambda_2\sqrt{\frac{\varphi^0-\inf F}{T\alpha}},
\]
by (B.62)--(B.63) with Lemma~B.11 in case (i) and by Lemma~B.12 in case (ii).
Here \(\lambda_1\) and \(\lambda_2\) depend only on \(L\), \(p\), \(\eta\),
\(\tau\), and the diameter \(2(L_h+1)\) of \(\mathcal Y\). Choose
\[
c:=\min\Bigl\{\frac1{36},\ \frac{(p-L)^2}{384p(p+L)^2},\
\Bigl(\frac{10}{11\lambda_1}\Bigr)^2\Bigr\},
\qquad
T:=\Bigl\lceil\Bigl(\frac{11}{10}\Bigr)^2\frac{\lambda_2^2\hat\Delta}{c\,\epsilon^4}\Bigr\rceil.
\]
Since \(\alpha=c\epsilon^2\le1\), case (i) gives
\(S_{t+1}\le\lambda_1\sqrt c\,\epsilon\le\frac{10}{11}\epsilon\), and case
(ii) gives \(S_{t+1}\le\frac{10}{11}\epsilon\) by Step~4 and the choice of
\(T\).

\paragraph{Step 6: from \(S_{t+1}\) to the KKT residual}
By~\eqref{eq:kkt-residual} and~\eqref{eq:prox-multiplier},
\(\mathcal R_{\mathrm{KKT}}(\bx^{t+1},\by^{t+1},\bq^t)
=\max\{\|\nabla_{\bx}\calD_\rho(\bx^{t+1},\blmbd^t)\|,\|\mathbf g\|\}\), where
\(\mathbf g:=\nabla_{\blmbd}\calD_\rho(\bx^{t+1},\blmbd^t)\) and
\(\blmbd^{t+1}-\blmbd^t=\tau\mathbf g\). By Step~2 and \(\tau/\rho\le1/144\),
\(\|\nabla_{\blmbd}\calD_\rho(\bx^{t+1},\blmbd^{t+1})-\mathbf g\|\le\|\mathbf g\|/144\),
hence \(\|\mathbf g\|\le\frac{144}{143}S_{t+1}\). Similarly,
\(G_{\calA}\tau\le1/144\) gives
\(\|\nabla_{\bx}\calD_\rho(\bx^{t+1},\blmbd^t)\|
\le S_{t+1}+G_{\calA}\tau\|\mathbf g\|\le\frac{144}{143}S_{t+1}\). Therefore
\(\mathcal R_{\mathrm{KKT}}(\bx^{t+1},\by^{t+1},\bq^t)\le\frac{11}{10}S_{t+1}\le\epsilon\),
and
\[
T_\epsilon(\theta,\alpha,1;I)\le t+1\le T\le C\epsilon^{-4},
\qquad
C:=\Bigl(\frac{11}{10}\Bigr)^2\frac{\lambda_2^2\hat\Delta}{c}+1,
\]
where the last inequality uses \(\epsilon\le1\). The constant \(c\) depends
only on \(L\), \(\gamma\), \(\eta\), \(\tau\), and \(L_h\), hence only on
\(\theta\) and \(\mathbf C\); the constant \(C\) depends in addition on
\(\Delta\) through \(\hat\Delta\). Neither depends on the instance or on the
dimensions.
\hfill\(\square\)

\subsection{Proof of Theorem~\ref{thm:upper-beta0}}
\label{app:upper-beta0}

All references to~\cite{li2026smoothing} in this appendix use the numbering
of arXiv:2602.14185v2. Let
\[
\ell:=\max\Bigl\{L_{\calD,\bx},G_{\calA},\frac1\rho\Bigr\},
\qquad r:=\frac1\gamma,\qquad \delta:=\frac1s,
\qquad \mathcal Y:=\{\blmbd\in\RR^m:\|\blmbd\|\le2L_h\},
\qquad D_{\mathcal Y}:=4L_h,
\]
where \(L_{\calD,\bx}\) is defined in~\eqref{eq:x-weak-convexity-modulus}.
We take \(L_h>0\). If the supplied class constant is zero, we use the
admissible Lipschitz bound \(L_h=1/16\) throughout this appendix; the
choice \(c_s=\max\{1,16L_h\}\) below remains \(1\), and the algorithm is
unchanged. All auxiliary constants are chosen from the fixed parameters
and class bounds, independently of the instance and its dimensions.

\paragraph{Step 1: reduction to Perturbed Smoothed GDA}
Since \(\beta=0\) and \(\bmu^0=\zero\), we have \(\bmu^k=\zero\) for all
\(k\). The potential function becomes
\[
\mathcal P_s(\bx,\blmbd;\bz,\zero)
=\calD_\rho(\bx,\blmbd)+\frac r2\|\bx-\bz\|^2
-\frac\delta2\|\blmbd\|^2.
\]
The correspondence with~\cite[(3.1) and Algorithm~3]{li2026smoothing} is
\[
f_{\mathrm{Li}}=\calD_\rho,\quad
(r_x,r_y^t,c^t,\alpha^t_{\mathrm{Li}},\beta^t_{\mathrm{Li}})
=(r,\delta,\eta,\tau,\alpha),\quad
\mathcal X=\RR^n.
\]
In particular, both algorithms use the old multiplier in the primal step,
the new primal iterate in the dual step, and the new primal iterate in the
center step. By~\eqref{eq:prox-multiplier}, the dual update reads
\[
\blmbd^{k+1}
=\Bigl(1-\frac\tau\rho-\frac\tau s\Bigr)\blmbd^k
+\frac\tau\rho\bq^k,
\qquad \|\bq^k\|\le L_h.
\]
The step-size conditions in Step~3 give
\(\tau\le1/(11(\ell+1/s))<1/(1/\rho+1/s)\).
Thus both coefficients are nonnegative and sum to \(1-\tau/s\le1\).
Induction yields \(\|\blmbd^k\|\le L_h\). Every multiplier lies in the
interior of \(\mathcal Y\), so the projection onto \(\mathcal Y\) does not
change the iterates and its normal cone at \(\blmbd^k\) is
\(\{\zero\}\). The primal projection is the identity. Consequently,
LRALM generates exactly the sequence of Perturbed Smoothed GDA applied to
\(\calD_\rho\) on \(\RR^n\times\mathcal Y\), from zero.

\paragraph{Step 2: assumptions and the unperturbed objective}
The gradient estimates in Step~2 of Appendix~\ref{app:upper-beta1} imply
\[
\begin{aligned}
\|\nabla_{\bx}\calD_\rho(\bx,\blmbd)
-\nabla_{\bx}\calD_\rho(\bx',\blmbd')\|
&\le\ell\bigl(\|\bx-\bx'\|+\|\blmbd-\blmbd'\|\bigr),\\
\|\nabla_{\blmbd}\calD_\rho(\bx,\blmbd)
-\nabla_{\blmbd}\calD_\rho(\bx',\blmbd')\|
&\le\ell\bigl(\|\bx-\bx'\|+\|\blmbd-\blmbd'\|\bigr).
\end{aligned}
\]
Hence Assumption~2.1 of~\cite{li2026smoothing} holds. Moreover,
\(\calD_\rho(\bx,\cdot)\) is concave and \(\mathcal Y\) is a compact
convex set containing zero, as required by its Assumption~2.2.
Taking \(\by=\calA(\bx)\) in~\eqref{eq:reduced-al} gives
\(\calD_\rho(\bx,\blmbd)\le F(\bx)\). For any
\(\blmbd\in\partial h(\calA(\bx))\), proximal optimality gives
\(\by(\bx,\blmbd)=\calA(\bx)\) and equality. Such a multiplier belongs
to \(\mathcal Y\), since its norm is at most \(L_h\). Therefore
\[
\max_{\blmbd\in\mathcal Y}\calD_\rho(\bx,\blmbd)=F(\bx),
\qquad \inf F>-\infty.
\]
The descent argument of~\cite[Section~4.2]{li2026smoothing} uses a lower
bound on the potential function. Accordingly, its minimum objective value
can be replaced by the infimum; attainment of \(\inf F\) is not needed.
The primal subproblems with quadratic regularization attain their minima
by strong convexity and coercivity.

\paragraph{Step 3: step-size conditions}
Set \(a_s:=r-\ell-\delta\) and \(b_s:=\ell+\delta\), and define
\[
\omega_s:=\frac1{\sqrt\delta}\,
\frac{a_s+\tau b_s(3r-2b_s)}{\tau a_s^{3/2}}.
\]
Under the correspondence in Step~1, Condition~3.2(ii)
of~\cite{li2026smoothing} reads
\begin{equation}
\begin{aligned}
&r=\Theta(\ell),\quad r>3\ell,\quad a_s>0,\quad \delta=1/s>0,\\
&\eta=\Theta(1/\ell),\quad 0<\eta<\frac1{r+\ell+\delta},\\
&\tau=\Theta(1/\ell),\quad
0<\tau\le\min\Bigl\{\frac1{11b_s},
\frac{\eta^2a_s^2}{4b_s(1+\eta a_s)^2}\Bigr\},\\
&\alpha=\Theta(\delta/\ell),\quad
0<\alpha\le\min\Bigl\{\frac1{36},
\frac{a_s^2}{384r(r+\ell+\delta)^2},
\frac1{384r\tau\omega_s^2}\Bigr\}.
\end{aligned}
\label{eq:beta0-stepsizes}
\end{equation}
The constants in these order relations are independent of \(s\).
For fixed \(r,\eta,\tau,\ell\) satisfying the first three conditions for
all sufficiently large \(s\), we have \(\omega_s^2=\Theta(s)\).
We can therefore choose a fixed, sufficiently small \(\varkappa>0\), as a
function of these fixed quantities, so that
\begin{equation}
\alpha:=\frac{\varkappa}{\ell s}
\label{eq:beta0-alpha}
\end{equation}
satisfies the last condition for all sufficiently large \(s\).
The conditions also give \(\eta<\gamma\) and \(\tau<s\), as required
by Definition~\ref{def:splalm}.

For example, all conditions hold for
\[
\gamma=\frac1{4\ell},\qquad \eta=\frac1{12\ell},\qquad
\tau=\frac1{800\ell},\qquad
\alpha=\frac1{2{,}000{,}000\,\ell s},\qquad
s\ge\max\{1,1/\ell\}.
\]
Indeed, \(2\ell\le a_s\le3\ell\) and \(b_s\le2\ell\), and
\[
\frac{\eta^2a_s^2}{4b_s(1+\eta a_s)^2}
\ge\frac1{392\ell}>\tau,\qquad
\frac{a_s^2}{384r(r+\ell+\delta)^2}\ge\frac1{13824\ell},
\qquad \omega_s^2\le800000\,\ell s.
\]
Thus \(1/(384r\tau\omega_s^2)\ge1/(1536000\,\ell s)>\alpha\),
and the remaining inequalities follow directly.

\paragraph{Step 4: initial gap for the perturbed objective}
Define
\[
\Phi_s(\bx):=\max_{\blmbd\in\mathcal Y}
\Bigl\{\calD_\rho(\bx,\blmbd)-\frac1{2s}\|\blmbd\|^2\Bigr\}.
\]
The maximization in~\eqref{eq:effective-penalty}, with \(\bmu=\zero\),
has an optimal multiplier \(s\bq/(\rho+s)\), where
\(\|\bq\|\le L_h\). This multiplier belongs to \(\mathcal Y\), so
restricting the maximization to \(\mathcal Y\) does not change its value.
It follows that
\begin{equation}
\Phi_s(\bx)=\calD_{\rho+s}(\bx,\zero),\qquad
F(\bx)-\frac{L_h^2}{2(\rho+s)}\le\Phi_s(\bx)\le F(\bx).
\label{eq:beta0-perturbed-objective}
\end{equation}
Let
\[
\begin{aligned}
v_0&:=\calD_\rho(\zero,\zero),\qquad
 d_0:=\min_{\bx}\Bigl\{\calD_\rho(\bx,\zero)
                         +\frac1{2\gamma}\|\bx\|^2\Bigr\},\\
p_0&:=\min_{\bx}\Bigl\{\Phi_s(\bx)+\frac1{2\gamma}\|\bx\|^2\Bigr\}.
\end{aligned}
\]
Since \(\mu_\gamma=1/\gamma-L_{\calD,\bx}>0\), the argument in
Step~4 of Appendix~\ref{app:upper-beta1} gives
\[
d_0\ge v_0-\frac{(G_f+G_{\calA}L_h)^2}{2\mu_\gamma},\qquad
p_0\le F(\zero),\qquad v_0\ge F(\zero)-\frac{L_h^2}{2\rho}.
\]
At zero initialization, the initial gap in
\cite[Definition~3.1(ii)]{li2026smoothing}, with the infimum convention
explained in Step~2, is
\(\Delta_{\Psi_2}=v_0-2d_0+2p_0-\inf\Phi_s\). Hence
\begin{equation}
\begin{aligned}
\Delta_{\Psi_2}
&\le\Delta+\frac{L_h^2}{2\rho}
+\frac{(G_f+G_{\calA}L_h)^2}{\mu_\gamma}
+\frac{L_h^2}{2(\rho+s)}\\
&\le\bar\Delta:=\Delta+\frac{L_h^2}{\rho}
+\frac{(G_f+G_{\calA}L_h)^2}{\mu_\gamma}.
\end{aligned}
\label{eq:beta0-initial-gap}
\end{equation}
The bound \(\bar\Delta\) is uniform over the problem class and independent
of \(s\), the target accuracy, and the dimensions.

\paragraph{Step 5: iteration complexity for game stationarity}
For \(k\ge0\), set
\[
S_{k+1}:=\max\Bigl\{
\|\nabla_{\bx}\calD_\rho(\bx^{k+1},\blmbd^{k+1})\|,
\|\nabla_{\blmbd}\calD_\rho(\bx^{k+1},\blmbd^{k+1})\|\Bigr\}.
\]
Both normal-cone terms vanish by Step~1, so this is the game-stationarity
measure for the unperturbed function \(\calD_\rho\) in
\cite[Theorem~3.3(i)]{li2026smoothing}.
The descent estimate in its Section~4.2 and estimates (4.19)--(4.23),
with~\eqref{eq:beta0-alpha}, give a constant
\(C_*=C_*(\eta\ell,\tau\ell,r/\ell,\varkappa)>0\) such that, among the first
\(K\ge1\) updates, some \(0\le k<K\) satisfies
\begin{equation}
S_{k+1}\le C_*\ell\sqrt{\frac{s\bar\Delta}{K}}
+\frac{D_{\mathcal Y}}s.
\label{eq:beta0-gs-rate}
\end{equation}
Here we take \(s\) sufficiently large that \(\ell s\ge1\), so the
term of order \(\sqrt{\ell\bar\Delta/K}\) is absorbed into the first
term. We average the descent estimate over \(k=0,\ldots,K-1\);
a fixed factor in this averaging is included in \(C_*\).
Thus the index in~\eqref{eq:beta0-gs-rate} refers to an output among the
first \(K\) updates of Definition~\ref{def:complexity-measure}.

Choose
\begin{equation}
c_s:=\max\{1,4D_{\mathcal Y}\}=\max\{1,16L_h\},\qquad
s:=\frac{c_s}{\epsilon},\qquad
K:=\max\Bigl\{1,\Bigl\lceil
16C_*^2\ell^2c_s\bar\Delta\,\epsilon^{-3}
\Bigr\rceil\Bigr\}.
\label{eq:beta0-budget}
\end{equation}
Then \(D_{\mathcal Y}/s\le\epsilon/4\), and the first term
in~\eqref{eq:beta0-gs-rate} is at most \(\epsilon/4\).
Consequently, some \(0\le k<K\) has \(S_{k+1}\le\epsilon/2\).

\paragraph{Step 6: from \(S_{k+1}\) to the KKT residual}
Write \(\mathbf d_k:=\blmbd^{k+1}-\blmbd^k\). The dual update and
Step~2 yield
\[
\|\mathbf d_k\|\le\tau\Bigl(S_{k+1}
+\frac{\|\mathbf d_k\|}\rho+\frac{L_h}s\Bigr),\qquad
\|\mathbf d_k\|\le\frac\tau{1-\tau/\rho}
\Bigl(S_{k+1}+\frac{L_h}s\Bigr).
\]
By~\eqref{eq:kkt-residual} and~\eqref{eq:prox-multiplier}, the output
residual is
\[
R_k:=\mathcal R_{\mathrm{KKT}}(\bx^{k+1},\by^{k+1},\bq^k)
=\max\Bigl\{\|\nabla_{\bx}\calD_\rho(\bx^{k+1},\blmbd^k)\|,
\|\nabla_{\blmbd}\calD_\rho(\bx^{k+1},\blmbd^k)\|\Bigr\}.
\]
Both gradients change by at most \(\ell\|\mathbf d_k\|\) when the
multiplier changes from \(\blmbd^{k+1}\) to \(\blmbd^k\). Since
\(\tau\ell\le1/11\) and \(\tau/\rho\le1/11\), we obtain
\begin{equation}
R_k\le\Bigl(1+\frac{\tau\ell}{1-\tau/\rho}\Bigr)S_{k+1}
+\frac{\tau\ell}{1-\tau/\rho}\frac{L_h}s
\le\frac{11}{10}S_{k+1}+\frac{L_h}{10s}.
\label{eq:beta0-residual-transfer}
\end{equation}
For the index selected in Step~5, \(c_s\ge16L_h\) gives
\[
R_k\le\frac{11}{20}\epsilon+\frac1{160}\epsilon
=\frac{89}{160}\epsilon<\epsilon.
\]
For all sufficiently small \(\epsilon\in(0,1]\), it follows that
\[
T_\epsilon(\theta(\epsilon),\alpha,0;I)
\le k+1\le K\le C'\epsilon^{-3},\qquad
C':=1+16C_*^2\ell^2c_s\bar\Delta.
\]
Finally, \(\alpha=\varkappa\epsilon/(\ell c_s)\), so
\(c_\alpha=\varkappa/(\ell c_s)>0\). Our fixed choice of \(\varkappa\)
ensures that \(c_\alpha\) depends only on
\(\rho,\gamma,\eta,\tau,\mathbf C\), and \(C'\) depends in addition
on \(\Delta\). Neither constant depends on the instance, the dimensions,
or \(\epsilon\).
\hfill\(\square\)

% Verification status (2026-10-14): the existence and parameter dependence
% of C_* in eq:beta0-gs-rate are taken from GPT Pro's conversion of Li26 v2,
% Section 4.2, (4.19)--(4.23), recorded in
% notes/gpt-proofs/beta0-upper-check-20261013.md, item 7.
% That conversion, including its fixed averaging factor, has not been
% independently verified line by line. Reading the cited statements and
% estimates does not constitute such a verification.
